% GENERATED FILE. Edit canonical lessons, then run npm run generate:books.
% canonical-source-hash: fnv1a64:d33e1e4a35a1dad3
% canonical-lessons: LA-C440-fluctuo, LA-C440-nauseo, LA-C440-defero, LA-C440-traicio, LA-C440-enato

\chapter{To toss on the waves, To be seasick, To carry down, To take across, To swim to safety}
\label{ch:la-to-toss-on-the-waves-to-be-seasick-to-carry-down-to-take-across-to-swim-to-safety}
\hlchaptermodality{440}

\begin{chapteropening}
\textbf{By the end of this chapter:} \emph{I can talk about a rough crossing: tossing on the waves, being seasick, being carried off course, crossing over, and swimming to safety.}

Complete the last lesson of chapter 440: I can talk about a rough crossing: tossing on the waves, being seasick, being carried off course, crossing over, and swimming to safety.
\end{chapteropening}

\section[flūctuō, flūctuāre]{flūctuō, flūctuāre — to toss on the waves, to rise in waves; to waver}

\label{lesson:LA-C440-fluctuo}



\begin{quote}
Before the new one: say the Latin for to launch, then the Latin for navigable.
\end{quote}

\subsection*{You'll want to know: flūctuō, flūctuāre}
\textbf{flūctuō, flūctuāre} — \textquotedblleft{}to toss on the waves, to rise in waves; to waver\textquotedblright{}.

\begin{quote}
In tempestāte nāvis diū flūctuat.
\end{quote}

\begin{quote}
\emph{In the storm the ship tosses on the waves for a long time.}
\end{quote}

\emph{flūctuō, flūctuāre, flūctuāvī}; also a deponent \emph{flūctuor}. Of the mind: to waver. English \emph{fluctuate} comes from it.

\begin{cousinweb}
\emph{flūctuō} is built on \emph{flūctus} 'wave' (EtymDB; Lewis \& Short), which comes from \emph{fluō} 'to flow'. English \emph{fluctuate} comes from it.
\end{cousinweb}

\begin{grammarlens}[title={What you've built}]
One more word for telling, explaining and describing at B1.
\end{grammarlens}

\subsection*{Your turn}
\begin{itemize}
\raggedright
  \item \emph{Say it:} \emph{flūctuō, flūctuāre}
  \item \emph{Say it:} \emph{flūctuō, flūctuāre}, once more
  \item \emph{Say it:} a sentence of your own with \emph{flūctuō, flūctuāre}
  \item \emph{Recall:} say the Latin for to launch, then the Latin for navigable, then say \emph{flūctuō, flūctuāre} again
\end{itemize}

\begin{tcolorbox}[breakable,colback=teal!4,colframe=teal!35!black,title={Before you move on}]
What does flūctuō, flūctuāre mean? (\textquotedblleft{}To toss on the waves, to rise in waves; to waver\textquotedblright{}.) Say it once more.
\end{tcolorbox}
\section[nauseō, nauseāre]{nauseō, nauseāre — to be seasick; to feel sick}

\label{lesson:LA-C440-nauseo}



\begin{quote}
Before the new one: say the Latin for navigable, then the Latin for to toss on the waves.
\end{quote}

\subsection*{You'll want to know: nauseō, nauseāre}
\textbf{nauseō, nauseāre} — \textquotedblleft{}to be seasick; to feel sick\textquotedblright{}.

\begin{quote}
Mare asperum est, et multī viātōrēs nauseant.
\end{quote}

\begin{quote}
\emph{The sea is rough, and many travellers are seasick.}
\end{quote}

\emph{nauseō, nauseāre, nauseāvī, nauseātum}, from \emph{nausea} 'seasickness'. English \emph{nauseate} comes from it.

\begin{cousinweb}
\emph{nauseō} is built on \emph{nausea} 'seasickness' (Lewis \& Short), which was borrowed from Greek \emph{nausiā}, built on \emph{naus} 'ship'. English \emph{nauseate} comes from it.
\end{cousinweb}

\begin{grammarlens}[title={What you've built}]
One more word for telling, explaining and describing at B1.
\end{grammarlens}

\subsection*{Your turn}
\begin{itemize}
\raggedright
  \item \emph{Say it:} \emph{nauseō, nauseāre}
  \item \emph{Say it:} \emph{nauseō, nauseāre}, once more
  \item \emph{Say it:} a sentence of your own with \emph{nauseō, nauseāre}
  \item \emph{Recall:} say the Latin for navigable, then the Latin for to toss on the waves, then say \emph{nauseō, nauseāre} again
\end{itemize}

\begin{tcolorbox}[breakable,colback=teal!4,colframe=teal!35!black,title={Before you move on}]
What does nauseō, nauseāre mean? (\textquotedblleft{}To be seasick; to feel sick\textquotedblright{}.) Say it once more.
\end{tcolorbox}
\section[dēferō, dēferre]{dēferō, dēferre — to carry down, to carry off; (of ships) to drive off course; to report}

\label{lesson:LA-C440-defero}



\begin{quote}
Before the new one: say the Latin for to toss on the waves, then the Latin for to be seasick.
\end{quote}

\subsection*{You'll want to know: dēferō, dēferre}
\textbf{dēferō, dēferre} — \textquotedblleft{}to carry down, to carry off; (of ships) to drive off course; to report\textquotedblright{}.

\begin{quote}
Ventus nāvem ad ignōtum lītus dēfert.
\end{quote}

\begin{quote}
\emph{The wind drives the ship off course to an unknown shore.}
\end{quote}

Irregular: \emph{dēferō, dēferre, dētulī, dēlātum}. In the passive, \emph{dēferrī}, of a ship carried off its course. English \emph{defer} comes from it.

\begin{cousinweb}
\emph{dēferō} is built from \emph{dē-} 'down, away' and \emph{ferō} 'to carry' (EtymDB; Lewis \& Short). English \emph{defer} comes from it through French.
\end{cousinweb}

\begin{grammarlens}[title={What you've built}]
One more word for telling, explaining and describing at B1.
\end{grammarlens}

\subsection*{Your turn}
\begin{itemize}
\raggedright
  \item \emph{Say it:} \emph{dēferō, dēferre}
  \item \emph{Say it:} \emph{dēferō, dēferre}, once more
  \item \emph{Say it:} a sentence of your own with \emph{dēferō, dēferre}
  \item \emph{Recall:} say the Latin for to toss on the waves, then the Latin for to be seasick, then say \emph{dēferō, dēferre} again
\end{itemize}

\begin{tcolorbox}[breakable,colback=teal!4,colframe=teal!35!black,title={Before you move on}]
What does dēferō, dēferre mean? (\textquotedblleft{}To carry down, to carry off; (of ships) to drive off course; to report\textquotedblright{}.) Say it once more.
\end{tcolorbox}
\section[trāiciō, trāicere]{trāiciō, trāicere — to take across, to cross over; to throw across}

\label{lesson:LA-C440-traicio}



\begin{quote}
Before the new one: say the Latin for to be seasick, then the Latin for to carry down.
\end{quote}

\subsection*{You'll want to know: trāiciō, trāicere}
\textbf{trāiciō, trāicere} — \textquotedblleft{}to take across, to cross over; to throw across\textquotedblright{}.

\begin{quote}
Mīlitēs equōs nāvibus trāns flūmen trāiciunt.
\end{quote}

\begin{quote}
\emph{The soldiers take their horses across the river in boats.}
\end{quote}

\emph{trāiciō, trāicere, trāiēcī, trāiectum}. \emph{Mare trāicere}, to cross the sea. English \emph{trajectory} comes from its participle.

\begin{cousinweb}
\emph{trāiciō} is built from \emph{trāns} 'across' and \emph{iaciō} 'to throw' (Lewis \& Short). English \emph{trajectory} comes from its participle \emph{trāiectus}.
\end{cousinweb}

\begin{grammarlens}[title={What you've built}]
One more word for telling, explaining and describing at B1.
\end{grammarlens}

\subsection*{Your turn}
\begin{itemize}
\raggedright
  \item \emph{Say it:} \emph{trāiciō, trāicere}
  \item \emph{Say it:} \emph{trāiciō, trāicere}, once more
  \item \emph{Say it:} a sentence of your own with \emph{trāiciō, trāicere}
  \item \emph{Recall:} say the Latin for to be seasick, then the Latin for to carry down, then say \emph{trāiciō, trāicere} again
\end{itemize}

\begin{tcolorbox}[breakable,colback=teal!4,colframe=teal!35!black,title={Before you move on}]
What does trāiciō, trāicere mean? (\textquotedblleft{}To take across, to cross over; to throw across\textquotedblright{}.) Say it once more.
\end{tcolorbox}
\section[ēnatō, ēnatāre]{ēnatō, ēnatāre — to swim to safety, to swim out}

\label{lesson:LA-C440-enato}



\begin{quote}
Before the new one: say the Latin for to carry down, then the Latin for to take across.
\end{quote}

\subsection*{You'll want to know: ēnatō, ēnatāre}
\textbf{ēnatō, ēnatāre} — \textquotedblleft{}to swim to safety, to swim out\textquotedblright{}.

\begin{quote}
Nāvis frangitur, sed nauta ad lītus ēnatat.
\end{quote}

\begin{quote}
\emph{The ship breaks up, but the sailor swims to the shore and is safe.}
\end{quote}

\emph{ēnatō, ēnatāre, ēnatāvī}. From \emph{ē-} 'out' and \emph{natō} 'to swim': to swim out of danger.

\begin{cousinweb}
\emph{ēnatō} is built from \emph{ē(x)-} 'out' and \emph{natō} 'to swim' (Lewis \& Short): to swim out of the water, to escape by swimming.
\end{cousinweb}

\begin{grammarlens}[title={What you've built}]
One more word for telling, explaining and describing at B1.
\end{grammarlens}

\subsection*{Your turn}
\begin{itemize}
\raggedright
  \item \emph{Say it:} \emph{ēnatō, ēnatāre}
  \item \emph{Say it:} \emph{ēnatō, ēnatāre}, once more
  \item \emph{Say it:} a sentence of your own with \emph{ēnatō, ēnatāre}
  \item \emph{Recall:} say the Latin for to carry down, then the Latin for to take across, then say \emph{ēnatō, ēnatāre} again
\end{itemize}

\begin{tcolorbox}[breakable,colback=teal!4,colframe=teal!35!black,title={Before you move on}]
What does ēnatō, ēnatāre mean? (\textquotedblleft{}To swim to safety, to swim out\textquotedblright{}.) Say it once more.
\end{tcolorbox}
